\documentclass[11pt,a4paper]{article}

\usepackage[margin=1in]{geometry}
\usepackage[T1]{fontenc}
\usepackage[utf8]{inputenc}
\usepackage{lmodern}
\usepackage{microtype}
\usepackage{amsmath,amssymb,mathtools}
\usepackage{booktabs}
\usepackage{graphicx}
\usepackage{float}
\usepackage[hidelinks]{hyperref}

\graphicspath{{../demos/report_outputs/exercise_1/}{../demos/report_outputs/exercise_2/}{../demos/report_outputs/exercise_3/}{Balint/demos/report_outputs/exercise_1/}{Balint/demos/report_outputs/exercise_2/}{Balint/demos/report_outputs/exercise_3/}}

\title{Final Report}
\author{}
\date{}

\begin{document}

\maketitle
\section*{Exercise 1: Entanglement Definitions and Probes}

    This section studies entanglement diagnostics in the periodic one-dimensional transverse-field Ising model using the project-owned exact-diagonalization workflow. The module structure is organized to reflect the same high-level API used in NetKet: graph, Hilbert space, operator, exact solver, sampler, and variational-state logic are kept conceptually separate so that later Monte Carlo workflows can reuse the same building blocks and so that a specific NQS workflow is simple to set up. For the one-dimensional TFIM, the expected asymptotic behavior is area-law saturation in the gapped phases and logarithmic block-entropy growth at criticality. The chosen field values are as follows: $h=0.5$ is the ordered reference point, $h=1.5$ is the disordered reference point, and $h=0.95$, $h=1.0$, and $h=1.05$ form the near-critical cluster used to describe the crossover region around the critical point.

    \subsection*{1/a Exact Diagonalization Setup}

    The exact-diagonalization path starts from the Hilbert-space definition provided by \texttt{SpinHilbert}. The current implementation is restricted to spin-$1/2$ systems, with configurations encoded as binary bitstrings with entries in $\{0,1\}$. The Hilbert-space object provides the total dimension $2^N$ and implements the map between basis states and integer indices. Site $0$ is stored as the least-significant bit, so the basis index of a configuration is
    \begin{equation}
        n = \sum_i s_i 2^i .
    \end{equation}
    For $L=16$, the Hilbert-space dimension is $2^{16}=65{,}536$, which is still small enough for exact diagonalization and explicit reshaping of the ground-state vector.

    The lattice geometry is supplied by the open \texttt{Chain1D} graph object, which provides a deterministic site ordering together with the nearest-neighbor bond list used in the Hamiltonian. This keeps the basis convention separate from the interaction pattern and makes the same infrastructure reusable for both one-dimensional chains and two-dimensional square lattices elsewhere in the project.

    The Hamiltonian is assembled through the TFIM operator builder, which constructs the operator as a sum of local $\sigma_i^z \sigma_j^z$ interaction terms over graph edges and on-site $\sigma_i^x$ field terms,
    \begin{equation}
        H = -J \sum_{\langle ij \rangle} \sigma_i^z \sigma_j^z - h \sum_i \sigma_i^x .
    \end{equation}
    The operator representation stores these contributions as local terms and exposes the connected matrix elements $H_{\sigma',\sigma}$ generated by acting on a basis configuration $\lvert \sigma \rangle$. This is important not only for exact diagonalization but also for the variational Monte Carlo estimator used later, where the local energy is written as
    \begin{equation}
        E_{\mathrm{loc}}(\sigma) = \sum_{\sigma'} H_{\sigma,\sigma'} \frac{\psi(\sigma')}{\psi(\sigma)} ,
    \end{equation}
    and the variational energy is obtained by averaging $E_{\mathrm{loc}}(\sigma)$ over configurations sampled from $\lvert \psi(\sigma) \rvert^2$.

    The exact solver then iterates over the operator matrix elements, builds a sparse COO matrix in coordinate-wise form, converts it to CSR format, and calls \path{scipy.sparse.linalg.eigsh} to obtain the lowest eigenpair using a Lanczos-based sparse eigensolver. The global phase of the ground state is fixed afterwards so that the amplitude with the largest magnitude is real. This sparse construction is preferable to forming a dense $65536 \times 65536$ matrix and already matches the connected-elements view that is reused later in the sampled VMC calculations.

    As a concrete example, the discussion below uses the half-chain split $A|B$ with $\lvert A \rvert = \lvert B \rvert = 8$. The Hilbert space then factorizes as
    \begin{equation}
        \mathcal{H} = \mathcal{H}_A \otimes \mathcal{H}_B ,
    \end{equation}
    with dimensions $2^{\lvert A \rvert}$ and $2^{\lvert B \rvert}$. Numerically, this means reshaping the exact state vector into a $2^{\lvert A \rvert} \times 2^{\lvert B \rvert}$ amplitude matrix $\psi(s_A,s_B)$, which is the key step for calculating reduced density matrices and Schmidt decompositions.

    \subsection*{1/b Von Neumann Entropy and the Reduced Density Matrix}

    For an exact pure state $\lvert \psi \rangle$, the reduced density matrix of subsystem $A$ is
    \begin{equation}
        \rho_A = \mathrm{tr}_B \lvert \psi \rangle \langle \psi \rvert ,
    \end{equation}
    and the von Neumann entropy is
    \begin{equation}
        S_1(A) = -\mathrm{tr}(\rho_A \log \rho_A) .
    \end{equation}
    The eigenvalues of $\rho_A$ are the squared Schmidt coefficients across the cut. When considering matrix product states, they quantify how much weight is carried by each Schmidt sector and therefore determine how compressible the state is across that bond. A rapidly decaying spectrum indicates that a small bond dimension would retain most of the state weight, while a broader spectrum implies weaker compressibility.

    The exact half-chain entropies for the retained field values are listed in Table~\ref{tab:exact-entropies}.

    \begin{table}[H]
        \centering
        \begin{tabular}{llrr}
            \toprule
            $h$  & Role                                 & $S_1(A)$ & $S_2(A)$ \\
            \midrule
            0.5  & ordered reference                    & 0.6990   & 0.6944   \\
            0.95 & near-critical point, ordered side    & 0.7762   & 0.6458   \\
            1.0  & critical point                       & 0.7501   & 0.5734   \\
            1.05 & near-critical point, disordered side & 0.6970   & 0.4842   \\
            1.5  & disordered reference                 & 0.3074   & 0.1419   \\
            \bottomrule
        \end{tabular}
        \caption{Exact half-chain entropies for the five retained TFIM field values on the open $L=16$ chain.}
        \label{tab:exact-entropies}
    \end{table}

    For the one-dimensional TFIM, the expected asymptotic distinction is that gapped phases obey area-law saturation while the critical point shows logarithmic growth with subsystem size. The present $L=16$ data do not cleanly expose that asymptotic separation: the half-chain entropy is largest at the near-critical ordered-side point $h=0.95$, while $h=1.0$ and $h=1.05$ remain numerically close to it. The useful interpretation is therefore that this finite open chain mainly resolves a crossover region around criticality rather than a clean thermodynamic-limit scaling law.

    The entanglement spectrum shown in Figure~\ref{fig:spectrum} supports the same interpretation.

    \begin{figure}[H]
        \centering
        \includegraphics[width=0.82\linewidth]{exercise_1_entanglement_spectrum.png}
        \caption{Leading Schmidt weights across the half-chain cut for the five retained field values.}
        \label{fig:spectrum}
    \end{figure}

    The spectrum becomes more concentrated as the field is increased toward the disordered regime. The Schmidt-compressibility data make this quantitative: keeping only the leading Schmidt value retains about $50.05\%$ of the weight at $h=0.5$, $64.74\%$ at $h=0.95$, $70.45\%$ at $h=1.0$, $75.91\%$ at $h=1.05$, and $93.02\%$ at $h=1.5$. The retained weight therefore increases away from the near-critical region, so truncation is easiest deep in the disordered phase and hardest on the ordered side and near criticality.

    \subsection*{1/c Renyi-2 Entropy and the SWAP Identity}
    $S_1$ remains primarily an exact-benchmark observable in the present context. Evaluating the von Neumann entropy requires explicit access to $\rho_A$ or all of its eigenvalues, and that becomes exponentially expensive with subsystem size. For small exact-diagonalization benchmarks it is straightforward; for a sampled neural quantum state it is not generally an efficient default observable, as it would reintroduce the same exponential wall that the variational ansatz is meant to avoid.

    The Renyi entropies are defined by
    \begin{equation}
        S_\alpha(A) = \frac{1}{1-\alpha} \log \mathrm{tr}(\rho_A^\alpha) .
    \end{equation}
    For $\alpha=2$, this reduces to
    \begin{equation}
        S_2(A) = -\log \mathrm{tr}(\rho_A^2) ,
    \end{equation}
    so the relevant quantity is the purity $\mathrm{tr}(\rho_A^2)$. The extremal limits are simple and useful checks: for a globally pure state that factorizes across the $A|B$ cut, $\rho_A$ is pure, so $\mathrm{tr}(\rho_A^2)=1$ and $S_2(A)=0$; for a maximally mixed reduced density matrix on a subsystem of dimension $d_A$, one has $\mathrm{tr}(\rho_A^2)=1/d_A$ and $S_2(A)=\log d_A$.

    The exact identity used later in the neural-quantum-state setting is
    \begin{equation}
        \mathrm{tr}(\rho_A^2)
        =
        \langle \psi \otimes \psi \rvert
        \mathrm{SWAP}_A
        \lvert \psi \otimes \psi \rangle .
    \end{equation}
    Writing the amplitudes as $\psi(s_A,s_B)$ gives
    \begin{equation}
        \rho_A(s_A,s_A') = \sum_{s_B} \psi(s_A,s_B)\psi^*(s_A',s_B) ,
    \end{equation}
    so that
    \begin{equation}
        \mathrm{tr}(\rho_A^2)
        =
        \sum_{s_A,s_A',s_B,s_B'}
        \psi(s_A,s_B)\psi^*(s_A',s_B)\psi(s_A',s_B')\psi^*(s_A,s_B') .
    \end{equation}
    This is exactly the replica expectation value of the operator that swaps the $A$ degrees of freedom between the two copies. Intuitively, the SWAP test asks how compatible two independent replicas remain after their subsystem-$A$ configurations are exchanged. If the state factorizes across the $A|B$ cut, swapping $A$ leaves the two-copy amplitude unchanged and the expectation is exactly $1$, so the subsystem is pure and the entropy ($log(1)$) is 0. If $A$ is strongly entangled with $B$, the swapped configurations are less compatible with the original amplitudes, so the expectation decreases, reflecting a more mixed reduced state on $A$ and increased entropy. The same construction can be generalized to $\alpha>2$ by introducing more replicas and a cyclic permutation. This is technically possible, but in the current implementation the estimator is already noisy enough that higher-$\alpha$ probes would not be practically useful even if an efficient estimator were written down.

    The SWAP-check data confirm the exact identity numerically for every retained field value. At $h=1.0$, for example, the purity from $\mathrm{tr}(\rho_A^2)$ is $0.5636148913$, while the purity reconstructed from $\exp(-S_2)$ is $0.5636148913$ to machine precision. The same agreement holds at $h=0.5$, $0.95$, $1.05$, and $1.5$, which validates the Renyi-2 route used later in the sampled neural-quantum-state setting.

    Because $S_2$ depends only on a low-order trace of $\rho_A$, it is much more practical for Monte Carlo estimation than $S_1$. One still expects the estimator to become harder for larger subsystems, but it avoids the need to reconstruct the full reduced density matrix or all of its eigenvalues.

    \subsection*{1/d Interpreting $S_1$ and $S_2$}

    For gapped one-dimensional ground states, entanglement is expected to obey an area law: for a fixed boundary cut, $S_1(A)$ and $S_2(A)$ should remain $\mathcal{O}(1)$ instead of growing linearly with subsystem size. A volume-law state would instead produce entropy proportional to $\lvert A \rvert$. Long-range entanglement is not simply ``large entropy''; it refers to nontrivial global entanglement structure that cannot be reduced to a short-range boundary contribution alone.

    The subsystem-size scan is shown in Figure~\ref{fig:scaling}.

    \begin{figure}[H]
        \centering
        \includegraphics[width=\linewidth]{exercise_1_entropy_scaling.png}
        \caption{Von Neumann and Renyi-2 entropy versus subsystem size for the five retained field values.}
        \label{fig:scaling}
    \end{figure}

    The main finite-size message from this plot is not a clean asymptotic scaling law, but a crossover comparison. The disordered reference point $h=1.5$ stays low across subsystem sizes and is the clearest area-law-like anchor in the retained data, which is what is expected for a gapped non-critical phase in the one-dimensional TFIM. At the same time, the finite system does not cleanly separate the close-to-critical points from the critical point itself: the near-critical trio $h=0.95$, $1.0$, and $1.05$ all show enhanced growth relative to $h=1.5$, with the largest finite-size entropies appearing slightly on the ordered side of criticality rather than exactly at $h=1.0$. The plot therefore serves mainly as a finite-size crossover window rather than as a clean asymptotic scaling test.

    The heuristic log-fit summary should be read cautiously. Its $R^2$ values are high for $h=0.95$, $1.0$, and $1.05$, which is consistent with enhanced finite-size growth around the crossover region, but the system is too small for strong claims about asymptotic critical scaling. The log-fit summary is therefore best read only as a compact finite-size diagnostic.

    Truncating the smallest Schmidt coefficients affects $S_1$ more strongly than $S_2$ because the two probes weight the entanglement spectrum differently. The von Neumann entropy contains the factor $-\lambda \log \lambda$, so even small Schmidt weights contribute. By contrast, $S_2$ depends on $\sum_i \lambda_i^2$, which emphasizes the dominant Schmidt sectors and suppresses the tail. This is why $S_1$ is the more tail-sensitive diagnostic, while $S_2$ is the more practical probe when moving to sampled neural quantum states.

    \subsection*{Conclusion}

    Exact diagonalization on the open $L=16$ TFIM chain provides a compact reference point for the entanglement probes used later in the project. The computational-basis convention, explicit graph structure, local-term Hamiltonian construction, and sparse Lanczos solve together define a consistent benchmark path that is also aligned with the connected-matrix-elements view required later in VMC.

    The reduced-density-matrix eigenvalues define the Schmidt spectrum across the half-chain cut and therefore quantify compressibility directly in matrix-product-state language. Across the retained field values, truncation becomes easier away from the critical regime, with the disordered point $h=1.5$ the most compressible case and the near-critical points remaining less compressible.

    The SWAP check verifies the Renyi-2 identity at machine precision for all five field values. This matters because $S_2$ remains a practical entanglement probe in the neural-quantum-state setting, whereas $S_1$ generally requires explicit reduced-density-matrix reconstruction. For the next exercises, this motivates carrying Renyi entropies into the neural-quantum-state setting.

    The finite-size scan is most naturally interpreted by taking $h=1.5$ as the low-entanglement area-law reference and $h=0.95$, $1.0$, and $1.05$ as the near-critical cluster for discussing crossover behavior. On this system size, however, the close-to-critical points and the critical point do not separate cleanly, and the largest half-chain entropies appear at $h=0.95$ rather than exactly at $h=1.0$. The contrast between $S_1$ and $S_2$ is nevertheless clear: $S_1$ is more sensitive to the tail of the entanglement spectrum, while $S_2$ is more robustly tied to the dominant Schmidt sectors.

    It's also worth keeping in mind that clearer scaling laws should emerge only at larger system sizes.

\section*{Exercise 2: Neural Quantum States and Sampled Renyi Entropy}

This section introduces the neural-quantum-state workflow used later in the project. The model module contains the variational ansaetze, the sampler module generates configurations with a local Metropolis chain, and the observables module evaluates quantities such as the SWAP-based Renyi entropy. This separation is a deliberate design choice: once a model returns a wavefunction amplitude for a spin configuration, the same sampling and measurement layer can be reused across all three architectures.

\subsection*{2/a Neural Quantum State Representation}

The common abstraction used throughout this exercise is a quantity of the form
\begin{equation}
    \log \psi_\theta(\sigma),
\end{equation}
which is the complex logarithm of the wavefunction amplitude for a spin configuration $\sigma$. This is the key interface boundary in the implementation. Once a model provides this object, the same sampling and observable code can be reused without model-specific logic in the Monte Carlo or SWAP-estimator layers.

All three ansatz classes first convert the stored spin variables from $\{0,1\}$ to the more standard $\{-1,+1\}$ convention. The final network output is then interpreted as two real channels: one for the log-amplitude and one for a raw phase variable $p_\theta(\sigma)$. In the implemented models the phase is bounded as $\phi_\theta(\sigma)=\pi \tanh p_\theta(\sigma)$, so the output takes the form
\begin{equation}
    \log \psi_\theta(\sigma) = a_\theta(\sigma) + i \pi \tanh p_\theta(\sigma).
\end{equation}
This keeps the interface uniform across all architectures and makes it straightforward to evaluate amplitude ratios later in the sampler and in the SWAP estimator.

For the present TFIM benchmark a real non-negative representation would in principle be sufficient in the computational basis, because the model is stoquastic there and its ground state can be chosen positive by a Perron--Frobenius argument.\footnote{See Bravyi et al., ``Monte Carlo simulation of stoquastic Hamiltonians,'' \emph{Quantum Information and Computation} 15 (2015), which explicitly lists the transverse-field Ising model as a stoquastic example.} The project still keeps a complex-valued interface because the same code is meant to carry over to frustrated models, where a positive real representation is not generally available. In particular, for the square-lattice $J_1$-$J_2$ antiferromagnet the Marshall--Peierls sign rule is known to break down once the frustration is strong enough, so the wavefunction cannot in general be treated as an everywhere-positive amplitude in the computational basis.\footnote{See Bishop, Farnell, and Parkinson, ``Phase Transitions in the Spin-Half $J_1$--$J_2$ Model,'' \emph{Phys. Rev. B} 58 (1998), which reports a breakdown of the Marshall--Peierls sign rule on the square lattice at $J_2/J_1 \approx 0.26$.}

The RBM ansatz is the closest to the standard NQS construction. Its hidden width is set by an integer density parameter $\alpha$, so the number of hidden units scales as $\alpha N$. After the visible spins are mapped to $\{-1,+1\}$, the hidden variables are summed out analytically, which gives the usual closed-form log-amplitude,
\begin{equation}
    \log \psi_\theta(\sigma)
    =
    \sum_i a_i \sigma_i
    +
    \sum_j \log \cosh\!\left(W_j \cdot \sigma + b_j\right)
    +
    i \pi \tanh \!\left[
        \sum_i \tilde{a}_i \sigma_i
        +
        \sum_j \left(\tilde{W}_j \cdot \sigma + \tilde{b}_j\right)
    \right].
\end{equation}
In the implemented RBM, the amplitude channel therefore keeps the standard analytically summed hidden contribution, while a separate linear phase channel supplies the complex phase before the final bounding step. Of the three models, this is the one that stays closest to the original RBM-NQS construction.

The FFNN ansatz replaces the analytic RBM structure by a dense multilayer perceptron. A configurable sequence of hidden widths defines the network depth and capacity, each hidden layer is followed by a fixed activation, and the final dense layer produces the two output channels for log-amplitude and phase. Unlike the RBM, there is no special hidden-spin summation here; the dependence on the spin configuration is learned entirely through the stacked nonlinear dense layers.

The CNN ansatz keeps the same output convention but changes the internal representation of the state. Instead of treating the spin configuration as a flat vector throughout, it is first reshaped into a lattice-shaped tensor, then passed through a stack of convolution layers, and finally flattened again before the two-channel output head. This is the only ansatz that explicitly takes the lattice geometry into account, so it can encode local geometric structure and translation-related patterns more naturally than a dense network.

\subsection*{2/b Architecture-Agnostic Monte Carlo Sampling}

Sampling is performed with a local Metropolis scheme acting directly on spin configurations. Starting from a batch of chains, one site is chosen uniformly at random in each chain and the corresponding spin is flipped, producing a proposed state $\sigma'$. The acceptance probability is computed from the wavefunction ratio,
\begin{equation}
    A(\sigma \rightarrow \sigma')
    =
    \min\!\left(
        1,
        \frac{\lvert \psi_\theta(\sigma') \rvert^2}{\lvert \psi_\theta(\sigma) \rvert^2}
    \right)
    =
    \min\!\left(
        1,
        \exp \bigl[ 2 \operatorname{Re} \bigl( \log \psi_\theta(\sigma') - \log \psi_\theta(\sigma) \bigr) \bigr]
    \right).
\end{equation}
This is exactly the expression used in the implementation: the log-amplitudes of the current and proposed states are evaluated in batch, and only their real parts enter the Metropolis ratio because the sampling measure is $\lvert \psi_\theta(\sigma) \rvert^2$.

The same sampler is used for RBM, FFNN, and CNN. In particular, no stochastic reconfiguration is introduced for the RBM in this exercise, for the sake of a unified VMC scheme for all three ansaetze.

The proposal move is a single-spin local update. This is the simplest option and the cleanest common baseline because it does not depend on any architecture-specific structure or model-specific proposal design. It also matches the project interface naturally: once $\log \psi_\theta(\sigma)$ is available, a local flip proposal can be scored without any extra model logic. In that sense it is the most transparent architecture-agnostic choice.

This simplicity comes with a cost. Alternatives such as multi-spin proposals, exchange moves, cluster-style updates, or more structured Hamiltonian-aware proposal rules could in principle improve mixing, especially when the state develops stronger correlations. A single-spin local move changes the configuration only minimally at each step, so one expects slower decorrelation when the correlation length becomes large. This is the usual critical-slowing-down problem of local Ising updates: near a continuous transition, increasingly large correlated regions have to be rearranged through many local moves.\footnote{This is the motivation behind cluster algorithms such as Swendsen--Wang; see Swendsen and Wang, ``Nonuniversal critical dynamics in Monte Carlo simulations,'' \emph{Phys. Rev. Lett.} 58 (1987).} Deep in the ordered phase the system is still strongly correlated, but it is no longer critical in the same sense, so the correlation length is finite and the slowdown is less severe than in the near-critical regime. In the present implementation, acceptance is therefore supplemented by a magnetization-autocorrelation diagnostic. The code does not yet estimate an integrated autocorrelation time, so acceptance alone should not be read as a full measure of sampler quality.

Figure~\ref{fig:sampler-acceptance} should be read in that limited sense. It shows that the sampler is numerically stable for the runs shown, but it is not by itself a full mixing study.

\begin{figure}[H]
    \centering
    \includegraphics[width=0.82\linewidth]{exercise_2_sampler_acceptance.png}
    \caption{Sampler acceptance for representative Exercise 2 runs.}
    \label{fig:sampler-acceptance}
\end{figure}

\subsection*{2/c Renyi-2 From the SWAP Estimator}

The practical entanglement observable carried over from Exercise 1 is the Renyi-2 entropy,
\begin{equation}
    S_2(A) = -\log \mathrm{tr}(\rho_A^2).
\end{equation}
For sampled neural quantum states, the useful identity is the replica-SWAP expression\footnote{For the original quantum Monte Carlo formulation of measuring Renyi entropy through a SWAP operator, see Hastings, Gonzalez, Kallin, and Melko, ``Measuring Renyi Entanglement Entropy in Quantum Monte Carlo Simulations,'' \emph{Phys. Rev. Lett.} 104 (2010).}
\begin{equation}
    \mathrm{tr}(\rho_A^2)
    =
    \langle \mathrm{SWAP}_A \rangle,
\end{equation}
which can be estimated from pairs of independently sampled configurations. Writing two configurations as $\sigma = (\sigma_A,\sigma_B)$ and $\eta = (\eta_A,\eta_B)$, the local SWAP estimator is
\begin{equation}
    \mathrm{SWAP}_A^{\mathrm{loc}}(\sigma,\eta)
    =
    \frac{
        \psi_\theta(\eta_A,\sigma_B)\,
        \psi_\theta(\sigma_A,\eta_B)
    }{
        \psi_\theta(\sigma_A,\sigma_B)\,
        \psi_\theta(\eta_A,\eta_B)
    }.
\end{equation}
Taking the average of this quantity over sampled replica pairs gives an estimator for $\langle \mathrm{SWAP}_A \rangle$, and therefore for $S_2(A)$.

The implementation follows this formula directly. Samples are flattened into a batch, paired row by row, and the subsystem-$A$ spins are exchanged between the two replicas to form the swapped configurations. The code then evaluates the log-amplitudes of the swapped states and constructs the estimator through exponentiated log-amplitude differences,
\begin{equation}
    \mathrm{SWAP}_A^{\mathrm{loc}}(\sigma,\eta)
    =
    \exp \Bigl[
        \log \psi_\theta(\eta_A,\sigma_B)
        +
        \log \psi_\theta(\sigma_A,\eta_B)
        -
        \log \psi_\theta(\sigma_A,\sigma_B)
        -
        \log \psi_\theta(\eta_A,\eta_B)
    \Bigr].
\end{equation}
This is exactly the form that makes the shared $\log \psi_\theta$ interface useful: the observable code only needs log-amplitude evaluations on original and swapped configurations, so no model-specific entanglement logic is required.

The small-system validation is summarized in Table~\ref{tab:swap-validation} and Figure~\ref{fig:swap-validation}. Here the exact $S_2(A)$ values come from explicit state reconstruction on the benchmark system, while the sampled $S_2(A)$ values are Monte Carlo estimates obtained from the SWAP estimator. The comparison therefore tests whether the sampled route reproduces the exact benchmark within the available noise level.

\begin{table}[H]
    \centering
    \begin{tabular}{rrrrr}
        \toprule
        $\lvert A \rvert$ & exact $S_2(A)$ & sampled $S_2(A)$ & sampled std & estimator std \\
        \midrule
        1 & 0.2246 & 0.2115 & 0.1188 & 0.0336 \\
        2 & 0.3212 & 0.3264 & 0.1901 & 0.0430 \\
        \bottomrule
    \end{tabular}
    \caption{SWAP-validation data against exact small-system benchmark values.}
    \label{tab:swap-validation}
\end{table}

\begin{figure}[H]
    \centering
    \includegraphics[width=0.82\linewidth]{exercise_2_swap_validation.png}
    \caption{Sampled SWAP estimates compared with exact Renyi-2 benchmark values.}
    \label{fig:swap-validation}
\end{figure}

The agreement is good at the level expected from the current Monte Carlo noise. For both retained subsystem sizes the sampled mean lies close to the exact value, with the remaining discrepancy at $\lvert A \rvert = 1$ still small compared with the observed sampling spread. This is enough to show that the implementation is targeting the correct quantity, while still making clear that the estimator is not quantitatively sharp without a larger sample budget.

That same point explains the noise discussion. As the subsystem grows, the swapped configurations become less typical under the original sampling distribution, so the ratio in the local estimator fluctuates more strongly. In practice this means that larger subsystems tend to produce noisier SWAP estimates. Higher-order Renyi entropies could be defined in the same replica spirit, but they would require more replicas and typically inherit even worse estimator quality. In the present setting, Renyi-2 is therefore the practical compromise between physical usefulness and estimator stability.

\subsection*{2/d Random Initialization Choices}

The first random-state study asks how strongly the sampled half-partition Renyi-2 depends on initialization before any training is performed. The two relevant choices are the overall parameter scale and whether the phase channel is allowed to fluctuate freely from the start. The summary is given in Table~\ref{tab:initialization-summary} and Figure~\ref{fig:initialization}.

\begin{table}[H]
    \centering
    \begin{tabular}{lrrrr}
        \toprule
        initialization & valid fraction & sampled $S_2(A)$ & sampled std & estimator std \\
        \midrule
        default complex & 0.00 & --- & 0.0000 & 0.0000 \\
        real scale 0.25 & 1.00 & 0.0409 & 0.0407 & 0.0415 \\
        real scale 1 & 1.00 & 1.4315 & 0.5413 & 0.8928 \\
        real scale 2 & 1.00 & 4.1525 & 1.5056 & 2.6525 \\
        \bottomrule
    \end{tabular}
    \caption{Summary of the random-initialization study.}
    \label{tab:initialization-summary}
\end{table}

\begin{figure}[H]
    \centering
    \includegraphics[width=0.9\linewidth]{exercise_2_initialization.png}
    \caption{Initialization comparison for the periodic $4\times 4$ TFIM random-state study.}
    \label{fig:initialization}
\end{figure}

The clearest conclusion is that unrestricted random complex phases are not usable under the reduced initialization-study budget. The default complex initialization produces no valid entropy points in the summary table, so allowing the phase channel to fluctuate freely from the outset pushes the SWAP estimator outside its numerically stable regime. This is consistent with the structure of the estimator: once the phase varies too strongly between original and swapped configurations, the replica ratio becomes much harder to average reliably.

Within the real-amplitude family, the overall scale matters strongly. A small real scale gives a state very close to a low-entanglement random baseline, while larger scales produce substantially larger sampled Renyi-2 values even before any optimization. In these data the half-partition estimate rises from about $0.041$ at scale $0.25$ to about $1.43$ at scale $1$ and about $4.15$ at scale $2$, although the uncertainty also grows sharply with scale. The interpretation is not that larger scales are automatically better, but that the random-state entanglement seen by the sampler is already highly sensitive to how broadly the amplitude channel is initialized.

\subsection*{2/e Architecture Comparison}

The second random-state study compares architecture families at a fixed real-amplitude initialization rule. The summary is listed in Table~\ref{tab:architecture-summary} and shown in Figure~\ref{fig:architecture}. Because the comparison is performed before optimization, it is best read as a comparison of representational and sampling behavior under a common initialization rule rather than as a definitive ranking of architectures.

\begin{table}[H]
    \centering
    \begin{tabular}{lrrr}
        \toprule
        model & sampled $S_2(A)$ & sampled std & estimator std \\
        \midrule
        CNN channels=4 & 0.0030 & 0.0021 & 0.0077 \\
        CNN channels=6 & 0.0037 & 0.0022 & 0.0093 \\
        FFNN width=8 & 0.0229 & 0.0129 & 0.0231 \\
        FFNN width=12 & 0.0320 & 0.0122 & 0.0320 \\
        RBM alpha=1 & 1.4854 & 0.2254 & 0.6003 \\
        RBM alpha=2 & 1.8217 & 0.3548 & 1.0675 \\
        \bottomrule
    \end{tabular}
    \caption{Summary of the architecture-comparison study.}
    \label{tab:architecture-summary}
\end{table}

\begin{figure}[H]
    \centering
    \includegraphics[width=0.9\linewidth]{exercise_2_architecture.png}
    \caption{Architecture comparison for the periodic $4\times 4$ TFIM random-state study.}
    \label{fig:architecture}
\end{figure}

The largest separation in the table is between the RBM and the other two architectures: the RBM produces order-one half-partition Renyi-2 values, whereas the FFNN and CNN remain close to zero. That observation is useful as a practical benchmark, but it should not be overinterpreted as a physics conclusion. Under the present initialization and sampling budget, the comparison mainly shows that the RBM explores substantially more entangled random states than the FFNN and CNN.

The within-family trends are modest for the FFNN and CNN families but clearer for the RBM. Increasing the FFNN width only changes the half-partition entropy slightly, and the two CNN choices remain close to the low-entanglement baseline. By contrast, the larger RBM reaches the highest sampled half-partition Renyi-2 in the comparison, which is why Exercise 3 carries forward the higher-capacity RBM variant.

\subsection*{Conclusion}

Exercise 2 sets up the neural-quantum-state side of the project around a common $\log \psi_\theta(\sigma)$ interface. That choice makes it possible to compare an RBM, a dense feed-forward network, and a convolutional network without changing the surrounding Monte Carlo and observable code. The local Metropolis sampler is kept deliberately simple and uniform across the three ansatz families, even though that simplicity comes with the expected weakness that single-spin updates decorrelate slowly as correlations grow.

The practical entanglement observable is the sampled Renyi-2 entropy obtained from the SWAP estimator. The small-system benchmark compares that sampled estimator directly against exact Renyi-2 values and shows that the method targets the correct observable, while also making clear that finite-sample noise is already a real issue even in a favorable setting. That is precisely why Renyi-2 is the natural quantity to carry forward: it remains accessible from samples, whereas more demanding entanglement probes would be even harder to estimate reliably.

The random-state studies make two additional points. First, initialization matters strongly, especially through parameter scale and through whether the phase channel is allowed to fluctuate freely. Second, complex amplitudes are not merely an implementation convenience: they are needed if the same framework is to extend beyond stoquastic benchmarks to frustrated models such as square-lattice $J_1$-$J_2$ systems, where a globally positive real representation is not generally available. Taken together, these results establish the measurement route used later in the project and clarify the practical constraints under which sampled entanglement estimates remain reliable.

\section*{Exercise 3: Variational Ground-State Search With RBM NQS}

This section studies variational ground-state optimization for the transverse-field Ising model within the neural-quantum-state workflow developed in the project. The calculation combines an RBM ansatz, local Metropolis sampling, the Adam optimizer, and Renyi-2 as the entanglement diagnostic tracked alongside the energy. The discussion is organized around the architecture choice, the VMC setup, a small-system benchmark where exact energies remain available as a clean reference, a larger-system training study, and two short outlook sections on GHZ targets and entanglement-spectrum tails.

\subsection*{3/a Architecture Choice}

The variational ansatz used in this section is \texttt{RBM(alpha=2)}. This choice is motivated by the architecture comparison from the previous exercise: under the same sampler and observable workflow, the RBM family was the only one that consistently produced clearly nontrivial half-partition Renyi-2 values.

Within that family, the higher-capacity choice \texttt{RBM(alpha=2)} is adopted here. The point is not that RBMs are universally optimal, but that they provide the most useful entanglement-carrying ansatz among the architectures tested in the same implementation.

The architecture summary supports the same reading numerically. In the comparison carried over from Exercise 2, the RBM gives the largest half-partition Renyi-2 among the three model families, while the FFNN and CNN remain substantially smaller.

\begin{figure}[H]
    \centering
    \includegraphics[width=0.82\linewidth]{exercise_3_architecture_summary.png}
    \caption{Architecture-comparison summary carried into the Exercise 3 ansatz choice.}
    \label{fig:architecture-summary}
\end{figure}

\subsection*{3/b VMC Setup For The TFIM Ground-State Search}

The variational objective is the ground-state energy of the neural quantum state,
\begin{equation}
    E_\theta = \frac{\langle \psi_\theta | H | \psi_\theta \rangle}{\langle \psi_\theta | \psi_\theta \rangle}.
\end{equation}
The code path follows the same module-based structure used in the earlier exercises. The model layer provides the RBM wavefunction through the common $\log \psi_\theta(\sigma)$ interface, the sampler module generates configurations with a local Metropolis chain, the observables module provides the Renyi-2 diagnostics, the optimizer module supplies the Adam update rule, and the VMC driver performs the step-by-step parameter updates.

The Hamiltonian expectation value is evaluated through the local-energy estimator. For each sampled configuration $\sigma$, the operator helper enumerates the connected configurations $\sigma'$ generated by the TFIM action, and the local energy is written as
\begin{equation}
    E_{\mathrm{loc}}(\sigma)
    =
    \sum_{\sigma'} H_{\sigma',\sigma}\,
    \frac{\psi_\theta(\sigma')}{\psi_\theta(\sigma)}.
\end{equation}
The gradient path is the standard log-derivative covariance estimator with
\begin{equation}
    O_\theta(\sigma) = \nabla_\theta \log \psi_\theta(\sigma),
\end{equation}
so that
\begin{equation}
    \nabla_\theta E_\theta
    =
    2\,\mathrm{Re}\!\left[
        \langle O_\theta^* E_{\mathrm{loc}} \rangle
        -
        \langle O_\theta^* \rangle
        \langle E_{\mathrm{loc}} \rangle
    \right].
\end{equation}
The VMC driver uses this sampled energy-and-gradient path directly, and Adam then turns the gradient estimate into a parameter update.

Sampling is performed with \texttt{MetropolisLocal} using 16 chains and 32 discard steps per chain. Energy is logged at every training step and the half-partition Renyi-2 is logged every two steps. The small benchmark uses a 256-step training budget with 256 measurement samples, while the larger sweep uses 70 training steps with 384 measurement samples.

One important point in this exercise is the status of the entanglement diagnostic. For the small proxy systems studied here, the trained NQS state can still be reconstructed exactly, so the subsystem-size scans use exact Renyi-2 of the trained variational state. That gives cleaner figures for these system sizes. The sampled SWAP route remains the relevant general method once exact reconstruction is no longer affordable.

\subsection*{3/c Small-System Benchmarks: Energy And Renyi-2 During Training}

The exact small-system benchmark uses a periodic $4\times 1$ TFIM chain and compares three field values under the same RBM and optimizer settings: a ferromagnetic point at $h = 0.5$, the critical point at $h = 1.0$, and a paramagnetic point at $h = 1.5$. This is the cleanest place to compare optimization behavior, because exact ground-state energies and exact half-partition Renyi-2 values are still available as references while the training histories remain easy to inspect directly. In one dimension, critical TFIM states are expected to show enhanced entanglement relative to gapped phases, with logarithmic scaling at criticality and area-law behavior away from it.\footnote{See Calabrese and Cardy, ``Entanglement Entropy and Quantum Field Theory,'' \emph{J. Stat. Mech.} P06002 (2004).}

\begin{table}[H]
    \centering
    \begin{tabular}{lrrrrr}
        \toprule
        run & final energy & exact energy & energy error & final $S_2$ & exact $S_2$ \\
        \midrule
        ferromagnetic $h=0.5$ & -4.269608 & -4.271558 & 0.001950 & 0.647272 & 0.659281 \\
        critical $h=1.0$ & -5.225886 & -5.226252 & 0.000366 & 0.386575 & 0.386461 \\
        paramagnetic $h=1.5$ & -6.758767 & -6.760009 & 0.001242 & 0.165577 & 0.167198 \\
        \bottomrule
    \end{tabular}
    \caption{Small-system summary for the periodic $4\times 1$ TFIM benchmark.}
    \label{tab:small-summary}
\end{table}

\begin{figure}[H]
    \centering
    \includegraphics[width=0.88\linewidth]{exercise_3_small_energy.png}
    \caption{Energy history for the periodic $4\times 1$ TFIM benchmark.}
    \label{fig:small-energy}
\end{figure}

\begin{figure}[H]
    \centering
    \includegraphics[width=0.88\linewidth]{exercise_3_small_entropy.png}
    \caption{Half-partition Renyi-2 history for the periodic $4\times 1$ TFIM benchmark.}
    \label{fig:small-entropy}
\end{figure}

With the longer training schedule, the small benchmark is now cleanly converged across all three field values. The final energy errors are of order $10^{-3}$ throughout, and the final half-partition Renyi-2 values also track the exact benchmarks closely. The critical run remains the most accurate of the three, but the ferromagnetic and paramagnetic runs now agree with the exact reference at the same qualitative level rather than lagging far behind it.

This improved benchmark gives a much clearer calibration of the retained RBM/VMC setup. On the periodic $4\times 1$ chain, the workflow is capable of reproducing both energies and coarse entanglement diagnostics to high accuracy once the optimization budget is made large enough. That result strengthens the interpretation of the later sampled-only sections: deviations there are much more plausibly due to large-system optimization and sampling difficulty than to a basic failure of the underlying ansatz or training loop.

\subsection*{3/d Larger-System Training}

The larger-system study keeps the same RBM, sampler, and optimizer, but shifts the emphasis from exact benchmarking to sampled training behavior. The retained systems are a periodic TFIM chain of length $32$, a periodic $5\times 5$ TFIM lattice, and a periodic $5\times 5$ square-lattice $J_1$-$J_2$ model. Exact energies and exact entanglement references are intentionally disabled in this section, so the results should be read as sampled-VMC benchmarks rather than as controlled exact comparisons.

The parameter choices are now aligned with the standard literature benchmarks. For the one-dimensional chain, the two retained points are the critical value $h=1.0$ and an off-critical paramagnetic point at $h=1.5$. For the square-lattice TFIM, the two-dimensional critical field is close to $h_c/J \approx 3.04$, so the notebook uses $h=3.0$ and $h=4.0$ to probe opposite sides of the transition.\footnote{For representative estimates of the square-lattice TFIM critical point, see for example Bl\"ote and Deng, ``Cluster Monte Carlo simulation of the transverse Ising model,'' \emph{Phys. Rev. E} 66 (2002), and subsequent tensor-network benchmarks reporting $h_c/J \approx 3.04$.} For the frustrated square-lattice $J_1$-$J_2$ model, the retained ratios are $J_1/J_2=1.8$ and $2.2$, corresponding to $J_2/J_1 \approx 0.56$ and $0.45$ and therefore straddling the commonly discussed strongly frustrated window near $J_2/J_1 \sim 0.5$.\footnote{For representative studies of the square-lattice spin-$1/2$ $J_1$-$J_2$ model and the intermediate nonmagnetic regime, see Hu, Becca, Parola, and Sorella, ``Direct evidence for a gapless $Z_2$ spin liquid by frustrating N\'eel antiferromagnetism,'' \emph{Phys. Rev. B} 88 (2013), and Bishop, Farnell, and Parkinson, ``Phase Transitions in the Spin-Half $J_1$--$J_2$ Model,'' \emph{Phys. Rev. B} 58 (1998).}

Figure~\ref{fig:large-graphs} summarizes the three lattice geometries used in the larger-system study. Blue edges denote nearest-neighbor couplings in all three cases, while the frustrated square-lattice model adds red diagonal second-neighbor bonds. For visual clarity, the two square-lattice graphs are drawn with open boundaries even though the simulated Hamiltonians use periodic boundary conditions. The one-dimensional chain is shown directly with periodic boundaries because the ring geometry remains visually clear at this size.

\begin{table}[H]
    \centering
    \begin{tabular}{lrrr}
        \toprule
        run & final energy & final $S_2$ & tail energy std. \\
        \midrule
        TFIM $L=32$, $h=1.0$ & -27.779286 & 1.266810 & 0.471005 \\
        TFIM $L=32$, $h=1.5$ & -13.777831 & 0.859072 & 0.740313 \\
        TFIM $5\times 5$, $h=3.0$ & -42.979151 & 0.221975 & 1.382138 \\
        TFIM $5\times 5$, $h=4.0$ & -49.251486 & 3.054304 & 4.485543 \\
        $J_1$-$J_2$ $5\times 5$, $J_1/J_2=1.8$ & -11.999737 & 0.979289 & 0.114004 \\
        $J_1$-$J_2$ $5\times 5$, $J_1/J_2=2.2$ & -27.194943 & 0.285725 & 0.165600 \\
        \bottomrule
    \end{tabular}
    \caption{Sampled larger-system summary for the periodic TFIM and $J_1$-$J_2$ benchmarks.}
    \label{tab:large-summary}
\end{table}

\begin{figure}[H]
    \centering
    \begin{minipage}[t]{0.32\linewidth}
        \centering
        \includegraphics[width=\linewidth]{exercise_3_tfim_chain_graph.png}
    \end{minipage}\hfill
    \begin{minipage}[t]{0.32\linewidth}
        \centering
        \includegraphics[width=\linewidth]{exercise_3_tfim_square_graph.png}
    \end{minipage}\hfill
    \begin{minipage}[t]{0.32\linewidth}
        \centering
        \includegraphics[width=\linewidth]{exercise_3_j1j2_square_graph.png}
    \end{minipage}
    \caption{Lattice graphs used to visualize the larger-system benchmarks. Left: periodic TFIM chain with $L=12$, drawn as a ring. Middle: $5\times 5$ TFIM lattice drawn with open boundaries for clarity, although the simulations use periodic boundaries. Right: $5\times 5$ $J_1$-$J_2$ lattice drawn with open boundaries for clarity; blue edges denote $J_1$ nearest-neighbor couplings and red edges denote $J_2$ diagonal couplings.}
    \label{fig:large-graphs}
\end{figure}

\begin{figure}[H]
    \centering
    \includegraphics[width=0.88\linewidth]{exercise_3_large_energy.png}
    \caption{Energy history across the periodic larger-system benchmark set.}
    \label{fig:large-energy}
\end{figure}

\begin{figure}[H]
    \centering
    \includegraphics[width=0.88\linewidth]{exercise_3_large_entropy.png}
    \caption{Half-partition Renyi-2 history across the periodic larger-system benchmark set.}
    \label{fig:large-entropy}
\end{figure}

\begin{figure}[H]
    \centering
    \includegraphics[width=0.9\linewidth]{exercise_3_large_entropy_scan.png}
    \caption{Post-training Renyi-2 versus subsystem size for the periodic larger-system benchmark set.}
    \label{fig:large-entropy-scan}
\end{figure}

The large-system figures are most useful as diagnostics of stability and qualitative entanglement ordering. In the one-dimensional TFIM chain, the expected ordering is preserved: the critical run at $h=1.0$ ends with a larger half-partition Renyi-2 than the off-critical run at $h=1.5$. This is the cleanest qualitative success in the sampled-only part of the exercise.

The two-dimensional runs are more uneven. On the periodic $5\times 5$ TFIM lattice, the $h=3.0$ run near the known critical region finishes with a modest sampled half-partition Renyi-2, while the $h=4.0$ run on the paramagnetic side ends with a much larger value and much larger late-stage energy fluctuations. The corresponding tail-energy standard deviations, $1.38$ and $4.49$, show that these two-dimensional TFIM runs are not equally well converged; the very large final $S_2$ at $h=4.0$ should therefore be interpreted as a warning sign of unstable sampled optimization rather than as a physically trustworthy entanglement ordering.

The frustrated $J_1$-$J_2$ pair is more stable than the two-dimensional TFIM pair. Both retained runs stay in a visibly narrower late-stage energy band, and their tail-energy standard deviations remain below $0.17$. Within this pair, the more frustrated point $J_1/J_2=1.8$ ends with the larger sampled half-partition Renyi-2 than $J_1/J_2=2.2$, which is qualitatively consistent with the expectation that stronger frustration near the intermediate regime enhances entanglement. Even so, these remain sampled diagnostics rather than controlled phase-boundary measurements.

\subsection*{3/e GHZ Diagnostic}

The GHZ section serves as a qualitative diagnostic of how the same RBM/VMC workflow behaves on a highly nonclassical target state, now phrased as the periodic $L=32$ TFIM chain at $h=0$.

On this large sampled-only problem the retained quantity of interest is the training history itself. The present workflow does not provide a reliable exact-state fidelity for this system size, and the sampled Renyi callback on this target still produces an unphysical negative entropy estimate in the summary output. The figure is therefore included only as evidence about optimization behavior on a cat-like target, not as a completed fidelity benchmark.

\begin{figure}[H]
    \centering
    \includegraphics[width=0.82\linewidth]{exercise_3_ghz_history.png}
    \caption{Energy history for the periodic $L=32$ TFIM chain at $h=0$, shown as a qualitative GHZ-style diagnostic.}
    \label{fig:ghz-history}
\end{figure}

\subsection*{3/f Entanglement-Spectrum-Tail Appendix}

The exact $4\times 1$ benchmark makes it possible to say something more precise about tail sensitivity than was possible in the larger sampled-only systems. Table~\ref{tab:tail-summary} and Figure~\ref{fig:tail-spectrum} compare the exact half-partition spectrum across the three small-system TFIM benchmarks and summarize both the von Neumann entropy $S_1$ and the Renyi-2 entropy $S_2$.

\begin{table}[H]
    \centering
    \begin{tabular}{lrrrr}
        \toprule
        run & exact $S_1$ & exact $S_2$ & $S_1-S_2$ & tail weight after top 2 \\
        \midrule
        TFIM $4\times 1$, $h=0.5$ & 0.681084 & 0.659281 & 0.021803 & 0.000734 \\
        TFIM $4\times 1$, $h=1.0$ & 0.521869 & 0.386461 & 0.135408 & 0.004278 \\
        TFIM $4\times 1$, $h=1.5$ & 0.309274 & 0.167198 & 0.142076 & 0.005566 \\
        \bottomrule
    \end{tabular}
    \caption{Exact tail-sensitive summary for the periodic $4\times 1$ TFIM benchmark.}
    \label{tab:tail-summary}
\end{table}

\begin{figure}[H]
    \centering
    \includegraphics[width=0.88\linewidth]{exercise_3_tail_spectrum.png}
    \caption{Exact half-partition entanglement spectrum for the periodic $4\times 1$ TFIM benchmark.}
    \label{fig:tail-spectrum}
\end{figure}

This comparison supports a conservative but useful conclusion. Exercise 2 already showed that sampled SWAP estimates reproduce exact small-system $S_2$ accurately enough for coarse entanglement benchmarking, and the present Exercise 3 benchmark retains that utility. However, $S_2$ is dominated by the largest reduced-density-matrix eigenvalues and is therefore only a blunt probe of the tail. The difference $S_1-S_2$ is much larger for the $h=1.0$ and $h=1.5$ ground states than for the $h=0.5$ state, even though all three are represented by the same four Schmidt levels on this cut. The tail weights after the two leading eigenvalues show the same pattern.

The practical implication is that agreement in energy and $S_2$ is necessary but not sufficient if the goal is to claim faithful reconstruction of the smallest Schmidt weights. The current workflow is therefore well supported as a coarse entanglement benchmark, but not yet as a tail-faithfulness benchmark for larger NQS states.

\subsection*{Conclusion}

Exercise 3 carries the project from static entanglement diagnostics to an actual variational ground-state search. The selected ansatz is \texttt{RBM(alpha=2)}, chosen because the RBM family already emerged as the most useful entanglement-capable model in the shared workflow. The VMC setup then combines that ansatz with the local Metropolis sampler, the connected local-energy estimator, the log-derivative gradient path, and the Adam optimizer.

The exact $4\times 1$ benchmark now shows that the retained RBM/VMC workflow can be quantitatively accurate across the ferromagnetic, critical, and paramagnetic points once the small-system optimization budget is sufficiently long. All three runs reach milliscale energy errors, and the final half-partition Renyi-2 values agree closely with the exact references. That benchmark therefore provides a strong calibration of the underlying ansatz and training loop before moving to larger sampled-only systems.

The larger-system study confirms that the workflow remains operational on a periodic $L=32$ TFIM chain, a periodic $5\times 5$ TFIM lattice, and a periodic $5\times 5$ $J_1$-$J_2$ model, but it also shows that interpretation must remain cautious without exact references. Some qualitative trends are physically sensible, especially the larger half-partition Renyi-2 in the critical one-dimensional TFIM run than in its off-critical counterpart and the stronger entanglement at $J_1/J_2=1.8$ than at $2.2$ in the frustrated square-lattice pair. By contrast, the $5\times 5$ TFIM comparison remains visibly unstable at the present budget, with the paramagnetic $h=4.0$ run carrying both the largest sampled Renyi-2 and the largest late-stage energy variance. The large-system figures are therefore best read as sampled benchmark diagnostics rather than definitive statements about phase-level entanglement ordering.

Finally, the tail appendix sharpens the entanglement interpretation. The current workflow is supported as a coarse benchmark through energy comparisons and reliable small-system $S_2$ checks, but the exact $4\times 1$ spectrum shows that $S_2$ alone does not determine the tail of the entanglement spectrum. That distinction is important for interpreting all later large-system NQS results: strong agreement in coarse observables is valuable, but it should not be overstated as proof of full spectral fidelity.
\end{document}
